\documentclass[runningheads]{llncs}

\usepackage[T1]{fontenc}
\usepackage{amsmath}
\usepackage{amsfonts}
\usepackage{booktabs}
\usepackage{graphicx}
\usepackage{multirow}
\usepackage{url}

\title{SASA-ESM: Surface-Aware Equivariant Residue Modeling for Protein--Protein Interaction Site Prediction}
\titlerunning{SASA-ESM for Protein Interface Prediction}

\author{Anonymous ICONIP Submission}
\authorrunning{Anonymous}
\institute{Paper under anonymous review}
% For the camera-ready version, replace the anonymous block with:
% \author{First Author\inst{1}\orcidID{0000-0000-0000-0000} \and Second Author\inst{1}}
% \authorrunning{F. Author et al.}
% \institute{Institution, City, Country \email{contact@example.com}}

\begin{document}
\maketitle

\begin{abstract}
Protein--protein interaction (PPI) site prediction requires residue-level reasoning over
biochemical exposure, sequence context, and three-dimensional geometry. We present
SASA-ESM, a reproducible pipeline that constructs weak residue labels from
\(\Delta\) solvent-accessible surface area (SASA), represents each residue with ESM-2
protein language model embeddings and local structural descriptors, and trains an
E(n)-equivariant graph neural network (EGNN) for interface classification. The pipeline
uses no external graph packages: SASA, \(\Delta\)SASA labels, residue features, graph
construction, EGNN training, optional cross-chain analysis, benchmark manifests, and prediction
exports are implemented in the accompanying code. On a 500-complex training corpus
(95,005 target-chain residues), the leakage-controlled ESM-2 650M plus geometry EGNN
with an 8 \AA{} residue graph cutoff reaches test F1 \(=0.7745\), AUROC \(=0.9322\),
and AUPRC \(=0.8421\). Adding apo SASA changes F1 only slightly to \(0.7787\).
In contrast, adding holo SASA raises F1 to \(0.8948\), demonstrating why holo-aware
features must be reported only as leakage diagnostics when labels are defined by
\(\mathrm{SASA}^{apo}-\mathrm{SASA}^{holo}\). Strict external evaluation reaches F1
\(=0.3529\) on Dset\_186-local and \(=0.3040\) on PDBtest\_315-local. These results
support a reproducible surface-labeled pipeline while exposing the remaining
generalization gap.

\keywords{Protein--protein interaction \and Solvent-accessible surface area \and
Protein language model \and ESM-2 \and Equivariant graph neural network}
\end{abstract}

\section{Introduction}

Identifying the amino-acid residues that form a protein--protein interface is a central
problem for structural bioinformatics, drug discovery, protein engineering, and functional
annotation. Experimental structures provide the most direct evidence of interacting
surfaces, but manually deriving residue-level labels is laborious and sensitive to the
definition of contact. A practical predictor must combine at least three complementary
signals: surface exposure, evolutionary sequence context, and geometric compatibility
between nearby residues.

Solvent-accessible surface area (SASA) is a natural physical measure for this task. A
residue that is exposed in an isolated chain and becomes buried in a complex is likely to
participate in the interface. This observation motivates the common \(\Delta\)SASA label:
\[
  \Delta\mathrm{SASA}_i =
  \mathrm{SASA}^{\mathrm{apo}}_i - \mathrm{SASA}^{\mathrm{holo}}_i,
  \quad
  y_i = \mathbb{1}[\Delta\mathrm{SASA}_i \ge \tau].
\]
However, SASA alone is not sufficient: exposed non-interface residues are frequent, and
interfaces depend on residue type, local secondary geometry, and spatial organization.
Recent protein language models such as ESM-2 encode rich sequence context
\cite{lin2023esm}, while equivariant graph neural networks can process residue
coordinates without discarding the symmetries of three-dimensional space
\cite{satorras2021egnn}.

This work integrates these ingredients into SASA-ESM. The contribution is deliberately
engineering-focused and reproducible: rather than assuming external preprocessing, the
project implements the full data path from PDB structures to residue-level predictions.
The main contributions are:
\begin{itemize}
  \item a weakly supervised \(\Delta\)SASA labeling pipeline for two-chain protein
  complexes;
  \item a leakage-controlled true ESM-2 650M residue representation combining
  1,280-dimensional protein language model embeddings, backbone/HSE descriptors,
  hydrophobicity, and C\(\alpha\) coordinates;
  \item an EGNN interface predictor, with cross-chain distance attention retained as
  an analysis module rather than a consistently improving core method;
  \item benchmark-ready prediction export on Dset\_186-local and PDBtest\_315-local,
  plus controlled feature and graph-cutoff ablations.
\end{itemize}

\section{Related Work}

Classical interface predictors often use geometric surface descriptors, conservation,
physicochemical residue properties, or hand-designed contact features. SASA-based
analysis traces back to the rolling-sphere definition of molecular surface accessibility
\cite{lee1971sasa}, and efficient implementations such as FreeSASA made large-scale
surface calculations practical \cite{mitternacht2016freesasa}. Graph-based deep learning
has since become a common choice for protein interface modeling because residue graphs
make local neighborhoods explicit. GraphPPIS is a representative residue graph model for
PPI site prediction and introduced widely used benchmark assets such as Test\_315
\cite{yuan2021graphppis}.

Protein language models add another axis of information. ESM-2 learns representations
from large protein sequence corpora and has been shown to encode structural and
functional regularities \cite{lin2023esm}. In parallel, equivariant networks such as
EGNN provide a principled way to use coordinates while respecting translation, rotation,
and reflection symmetries \cite{satorras2021egnn}. SASA-ESM combines these lines: the
label is defined by a physical surface criterion, the node representation includes
language-model and local structural features, and the graph model is equivariant.

\section{Method}

\subsection{Data Construction and Labeling}

The corpus is built from two-chain protein assemblies downloaded from the RCSB Protein
Data Bank \cite{berman2000pdb}. Each complex is parsed into atoms, residues, chains, and
C\(\alpha\) residue coordinates. We retain target and partner chains with at least 30
residues to avoid degenerate small fragments.

For each target-chain residue, SASA is calculated twice. The apo value is computed on the
target chain alone. The holo value is computed after placing the target chain in the
presence of the partner chain. A residue is labeled as interface when
\(\Delta\mathrm{SASA}\ge 2.0\) \AA\({}^2\). This threshold is used throughout the
reported experiments. Importantly, the strict protocol uses SASA only to define weak
labels. Neither \(\Delta\)SASA nor its holo component is used as an input feature. We
separately report apo-only and holo-aware runs to quantify label-definition leakage.

\subsection{Residue Representation}

In the strict primary model, each residue node is represented by:
\[
  x_i = [e_i,\,s_i],
\]
where \(e_i\in\mathbb{R}^{1280}\) is the ESM-2 650M residue embedding and \(s_i\)
contains backbone dihedral encodings \((\sin\phi,\cos\phi,\sin\psi,\cos\psi)\), half
sphere exposure counts, and hydrophobicity. The resulting strict true-650M feature set
contains 1,287 scalar input features per residue. An apo-only ablation appends
\(\mathrm{SASA}^{apo}_i\), while the explicitly labeled holo-aware diagnostic appends
both apo and holo SASA. The code validates dimensionality to avoid
the common failure mode of treating smaller 320-dimensional ESM features as 650M
features.

\subsection{Equivariant Residue Graph}

For each target chain, a residue graph \(G=(V,E)\) is built from C\(\alpha\)
coordinates. Residues \(i\) and \(j\) are connected when
\(\|p_i-p_j\|_2 \le r\), where \(r\) is the distance cutoff. We evaluate
\(r\in\{6,8,10,12\}\) \AA{}.

The EGNN updates node states and coordinates through messages
\[
  m_{ij} = \phi_e(h_i, h_j, \|p_i-p_j\|_2/10),
\]
with residual feature updates and bounded coordinate weights. The implementation uses
unit direction vectors and normalized distances for numerical stability. After three
EGNN layers, an MLP classifier outputs an interface logit for every residue.

\subsection{Cross-Chain Spatial Attention}

The cross-chain variant augments the target-chain EGNN output with partner-chain C\(\alpha\)
positions. For target residue \(i\) and partner residue \(j\), attention is computed as
\[
  \alpha_{ij} =
  \mathrm{softmax}_j(-\|p_i-q_j\|_2/\sigma),
\]
where \(\sigma\) is learnable. Partner coordinates are projected into hidden vectors and
combined with target hidden states before classification. This lightweight mechanism
lets the model access the spatial layout of the interacting chain without constructing a
full bipartite atom graph.

\begin{figure}[t]
\centering
\fbox{\begin{minipage}{0.93\linewidth}
\centering
\small
PDB complex
\(\rightarrow\) apo/holo SASA
\(\rightarrow\) \(\Delta\)SASA labels
\(\rightarrow\) ESM-2 650M residue embeddings
\(\rightarrow\) ESM + structural node features
\(\rightarrow\) EGNN
\(\rightarrow\) residue interface probabilities
\end{minipage}}
\caption{SASA-ESM processing pipeline. The same manifest drives label generation,
embedding extraction, multimodal dataset construction, training, and benchmark
prediction export.}
\label{fig:pipeline}
\end{figure}

\section{Experimental Setup}

\subsection{Datasets}

Table~\ref{tab:datasets} summarizes the datasets used in this manuscript. The main
corpus contains 500 two-chain complexes and 95,005 labeled target-chain residues. The
Dset\_186 evaluation uses the locally analyzable subset produced by the benchmark
pipeline after structure availability and chain-quality filtering: 158 complexes and
39,505 residues. PDBtest\_315-local uses 314 analyzable complexes and 65,119 residues
after the same chain-quality rules. We use the \emph{-local} suffix because both datasets
are produced by this repository's reproducible preprocessing protocol.

\begin{table}[t]
\centering
\caption{Datasets used in the current package.}
\label{tab:datasets}
\begin{tabular}{lrrl}
\toprule
Dataset & Complexes & Residues & Role \\
\midrule
Main two-chain corpus & 500 & 95,005 & Train/validation/test split \\
Dset\_186-local & 158 & 39,505 & External benchmark prediction \\
PDBtest\_315-local & 314 & 65,119 & External benchmark prediction \\
\bottomrule
\end{tabular}
\end{table}

\subsection{Training Protocol}

Models are trained with Adam, learning rate \(10^{-3}\), weight decay \(10^{-4}\),
hidden dimension 128, dropout 0.4, three EGNN layers, gradient clipping at 1.0, and
early stopping based on validation F1. Complexes are split by complex identity into
70/15/15 train/validation/test partitions with seed 42. Metrics include accuracy,
precision, recall, F1, AUROC, and AUPRC. AUPRC is reported because interface labels are
class-imbalanced and precision--recall curves are more informative under skewed class
priors \cite{davis2006pr}.

\section{Results}

\subsection{Leakage-Controlled Main Corpus}

Table~\ref{tab:main_results} reports strict baselines and leakage diagnostics. The
strict ESM plus geometry EGNN improves over MLP and GCN baselines. Adding apo SASA has
only a small effect. Adding holo SASA produces a much larger score increase because the
label itself is computed from apo minus holo SASA; this configuration is diagnostic and
is not treated as the deployable result.

\begin{table}[t]
\centering
\caption{Held-out test performance. Holo-aware rows are leakage diagnostics, not strict predictors.}
\label{tab:main_results}
\resizebox{\textwidth}{!}{%
\begin{tabular}{llrrrr}
\toprule
Role & Model / features & F1 & AUROC & AUPRC & Inputs \\
\midrule
Strict baseline & apo-SASA rank & 0.3317 & 0.5841 & 0.2550 & apo SASA \\
Strict baseline & MLP / apo & 0.4292 & 0.7000 & 0.3098 & apo SASA \\
Strict baseline & MLP / ESM & 0.6290 & 0.8949 & 0.7387 & ESM \\
Strict baseline & GCN / ESM+geom. & 0.5961 & 0.8730 & 0.6570 & ESM+geom. \\
Strict primary & EGNN / ESM+geom. & \textbf{0.7745} & \textbf{0.9322} & \textbf{0.8421} & ESM+geom. \\
Strict ablation & EGNN / apo+ESM+geom. & 0.7787 & 0.9348 & 0.8514 & apo+ESM+geom. \\
Diagnostic & MLP / apo+holo & 0.8724 & 0.9457 & 0.9117 & apo+holo SASA \\
Diagnostic & EGNN / apo+holo+ESM+geom. & 0.8948 & 0.9768 & 0.9497 & holo-aware \\
\bottomrule
\end{tabular}
}
\end{table}

\subsection{Distance Cutoff Ablation}

The graph cutoff controls the amount of geometric context available to each residue.
Table~\ref{tab:cutoff} shows that 8 \AA{} yields the best F1 and AUPRC, while larger
cutoffs do not improve the held-out test set. This suggests that the most useful
information is local but not extremely short-range; a 6 \AA{} graph may miss relevant
neighbors, whereas 10--12 \AA{} adds edges with weaker interface specificity.

\begin{table}[t]
\centering
\caption{Holo-aware EGNN cutoff diagnostic. These rows are not strict predictors.}
\label{tab:cutoff}
\begin{tabular}{rrrrrr}
\toprule
Cutoff & Acc. & Prec. & Rec. & F1 & AUROC / AUPRC \\
\midrule
6 \AA{}  & 0.9574 & 0.9027 & 0.8781 & 0.8902 & 0.9750 / 0.9420 \\
8 \AA{}  & \textbf{0.9590} & 0.9023 & \textbf{0.8874} & \textbf{0.8948} & \textbf{0.9768} / \textbf{0.9497} \\
10 \AA{} & 0.9577 & 0.9040 & 0.8777 & 0.8907 & 0.9758 / 0.9488 \\
12 \AA{} & 0.9584 & \textbf{0.9188} & 0.8645 & 0.8908 & 0.9727 / 0.9449 \\
\bottomrule
\end{tabular}
\end{table}

\subsection{Dset\_186-local Prediction}

Table~\ref{tab:dset186} reports strict prediction metrics on Dset\_186-local. The
holo-aware rows are retained only to show label-definition leakage and to analyze the
optional cross-chain module. The strict drop exposes a substantial generalization gap
that must be addressed before claiming competitive benchmark performance.

\begin{table}[t]
\centering
\caption{External prediction on Dset\_186-local (158 complexes, 39,505 residues).}
\label{tab:dset186}
\resizebox{\textwidth}{!}{%
\begin{tabular}{llrrrr}
\toprule
Role & Model & Prec. & Rec. & F1 & AUROC / AUPRC \\
\midrule
Strict primary & EGNN & 0.3605 & 0.3456 & 0.3529 & 0.7277 / 0.3050 \\
Diagnostic & Holo-aware EGNN & 0.6690 & 0.6223 & 0.6448 & 0.8868 / 0.6770 \\
Diagnostic & Holo-aware cross-chain EGNN & 0.6167 & 0.6735 & 0.6439 & 0.8891 / 0.5766 \\
\bottomrule
\end{tabular}
}
\end{table}

\subsection{PDBtest\_315-local Prediction}

Table~\ref{tab:pdbtest315} reports the completed PDBtest\_315-local run. All 315 official
chain-aware entries obtained structures from RCSB; 314 passed the current chain-quality
rules. ESM-2 650M embeddings and both model predictions were produced on one NVIDIA
GeForce RTX 5060 Laptop GPU. The strict model is reported first. Holo-aware EGNN and
cross-chain rows are leakage diagnostics; within that diagnostic setting, cross-chain
attention trades precision and AUPRC for recall and F1.

\begin{table}[t]
\centering
\caption{External prediction on PDBtest\_315-local (314 complexes, 65,119 residues).}
\label{tab:pdbtest315}
\resizebox{\textwidth}{!}{%
\begin{tabular}{llrrrr}
\toprule
Role & Model & Prec. & Rec. & F1 & AUROC / AUPRC \\
\midrule
Strict primary & EGNN & 0.3057 & 0.3023 & 0.3040 & 0.6799 / 0.2675 \\
Diagnostic & Holo-aware EGNN & 0.7360 & 0.6252 & 0.6761 & 0.8946 / 0.7408 \\
Diagnostic & Holo-aware cross-chain EGNN & 0.7014 & 0.6629 & 0.6816 & 0.8983 / 0.6906 \\
\bottomrule
\end{tabular}
}
\end{table}

\section{Discussion}

The experiments show that a physically grounded label source and a language-model
representation can form a reproducible engineering pipeline, but they also expose a
critical evaluation hazard. Holo SASA is too close to the \(\Delta\)SASA label
definition to be accepted as a deployable input. ESM-2 contributes sequence and
evolutionary context; structural descriptors and coordinates let the EGNN model local
geometry. The true-650M correction is important: the final strict files contain 1,280
ESM dimensions and 1,287 total features,
whereas older 320-dimensional embeddings would correspond to a different ESM-2 scale and
cannot be compared as 650M experiments.

Cross-chain attention is therefore treated as an analysis module, not a stable core
improvement. In the holo-aware diagnostic it improves recall and F1 on
PDBtest\_315-local, but does not consistently improve precision or AUPRC. A future
strict training setup should explicitly sample non-interacting partner candidates.

\section{Limitations and Reproducibility}

This manuscript reports only experiments completed in the packaged project state.
PDBtest\_315-local preprocessing, ESM-2 extraction, and prediction are complete. The
collector can be configured for 1,000+ complexes and the extractor supports larger ESM-2
models such as 3B, but those scale-up experiments should be run and reported separately
with storage, runtime, and hardware details.

GraphPPIS is discussed as a representative prior method, but its published values are
not inserted into our numeric tables because the current local preprocessing and split
protocols are not guaranteed to match its official evaluation protocol. A direct
comparison requires running both methods on the same manifests and labels.

All reported numeric tables are backed by CSV files included with this submission
package under \texttt{data/}. The implementation exports per-residue predictions for
benchmark analysis and stores model checkpoints with feature columns, normalization
statistics, cutoff, hidden size, and metric metadata.

\section{Conclusion}

SASA-ESM provides an end-to-end implementation for weakly supervised PPI site
prediction from protein complex structures. Under strict no-SASA inputs, true ESM-2
650M features and EGNN geometry reach internal F1 \(=0.7745\). External results are
substantially lower and should be read as a reproducible baseline rather than a
state-of-the-art claim. The leakage audit, exported predictions, and explicit diagnostic
rows make the current package a more defensible foundation for future work.

\bibliographystyle{splncs04}
\bibliography{references}

\end{document}
